% main_zh.tex -- 中文稿件（与 main.tex 同步的双根之一；ctexart，A4 单栏 11pt，行距 1.5）。编译：latexmk -xelatex main_zh.tex
% 正文在 sections/zh/，与 sections/en/ 文件一一对应；表格、宏、图、附录术语表与英文稿共用同一份文件。
\documentclass[11pt,a4paper,fontset=fandol]{ctexart}
\newif\iflangzh\langzhtrue
\newcommand{\bi}[2]{#1}            % 共享文件写 \bi{中文}{English}，本根取中文
\usepackage[a4paper,top=2.5cm,bottom=2.5cm,left=2.6cm,right=2.6cm]{geometry}
\usepackage{fontspec}
\IfFontExistsTF{Times New Roman}{\setmainfont{Times New Roman}}{\setmainfont{TeX Gyre Termes}}
\usepackage{amsmath,amssymb,amsthm,bm}
\usepackage{booktabs,multirow,array,tabularx,threeparttable}
\usepackage{graphicx}
\usepackage{enumitem}
\setlist{itemsep=2pt,topsep=4pt,parsep=0pt}
\usepackage{setspace}
\setstretch{1.5}
\usepackage{float}
\usepackage{caption}
\captionsetup[table]{name=表,font=small,labelfont=bf,skip=6pt,position=top}
\captionsetup[figure]{name=图,font=small,labelfont=bf,skip=6pt,position=bottom}
\usepackage{algorithm}
\usepackage{algpseudocode}
\floatname{algorithm}{算法}
\renewcommand{\algorithmicrequire}{\textbf{输入：}}
\renewcommand{\algorithmicensure}{\textbf{输出：}}
\usepackage[colorlinks=true,linkcolor=deepblue,citecolor=deepblue,urlcolor=deepblue]{hyperref}
\newcolumntype{L}{>{\raggedright\arraybackslash}X}
\theoremstyle{plain}
\newtheorem{theorem}{定理}
\newtheorem{proposition}{命题}
\newtheorem{conjecture}{猜想}
\renewcommand{\proofname}{证明}
\newif\ifanonymous
% macros/placeholders.tex -- placeholder and draft-stage macros shared by main.tex and main_zh.tex.
% \bi{中文}{English} is defined by each root before this file is read: the Chinese root keeps the
% first argument, the English root the second. Every shared file (tables/, appendix/, this file)
% writes its visible text through \bi so that one source serves both manuscripts.
\usepackage{xcolor}
\definecolor{deepblue}{RGB}{0,64,128}
\definecolor{phgray}{RGB}{120,120,120}
\definecolor{warnbg}{RGB}{255,247,230}
\definecolor{citered}{RGB}{160,30,30}
% numeric placeholders: single value / mean +- std / expected-best (bold) slot
\newcommand{\ph}{\textcolor{phgray}{--}}
\newcommand{\phpm}{\textcolor{phgray}{$\text{--}\!\pm\!\text{--}$}}
\newcommand{\phb}{\textcolor{phgray}{\textbf{--}}}
% a claim that no experiment has judged yet
\newcommand{\pending}[1]{#1\textsuperscript{\textcolor{deepblue}{[\bi{待裁决}{pending}]}}}
% a citation whose entry has not been verified: never a fabricated bib key
\newcommand{\citeneed}[1]{\textcolor{citered}{[\bi{待核实引文}{ref}: #1]}}
% a boxed note that the draft stage must carry (title page, contribution list, experiments)
\newcommand{\draftnote}[1]{%
  \begin{center}\setlength{\fboxsep}{6pt}%
  \colorbox{warnbg}{\parbox{0.93\linewidth}{\small\textbf{\bi{草稿提示：}{Draft note: }}#1}}%
  \end{center}}
% figure placeholder = the drawing brief handed to opt-paper-scifig
% {content elements}{data source}{route and design width}{box height}
\newcommand{\figplace}[4]{%
  \begin{center}\setlength{\fboxsep}{6pt}%
  \fbox{\parbox[c][#4][c]{0.9\linewidth}{%
    \centering\small
    {\bfseries\bi{【图形占位·待 opt-paper-scifig 生成】}{[Figure placeholder -- to be produced by opt-paper-scifig]}}\\[4pt]
    \begin{minipage}{0.95\linewidth}\footnotesize\raggedright
      \textbf{\bi{内容要素：}{Elements: }}#1\par
      \textbf{\bi{数据来源：}{Data: }}#2\par
      \textbf{\bi{生成路线：}{Route: }}#3
    \end{minipage}}}%
  \end{center}}

% generated by paper_assets/scripts/make_tables.py -- do not edit by hand

\newcommand{\OursGapMain}{\ph{}}
\newcommand{\BestBaselineGapMain}{\ph{}}
\newcommand{\ImpvOverBest}{\ph{}}
\newcommand{\OursDecodeReduction}{\ph{}}
\newcommand{\NumFixedInstances}{59}
\newcommand{\NumRuns}{3}

% frontmatter-authors.tex -- generated by paper_assets/scripts/make_submission.py from submission.yaml.
% Placeholder (anonymous) until the submission package is produced. No names, e-mails or ORCIDs here.
\global\anonymoustrue
\newcommand{\papercredit}{}
\newcommand{\paperrepo}{}
\newcommand{\paperfunding}{}

\begin{document}
\begin{center}
{\LARGE\bfseries 面向多重度柔性作业车间调度的流体-残差进化搜索\par}
\vspace{6pt}
{\large Fluid-residual evolutionary search for multiplicity flexible job-shop scheduling\par}
\vspace{10pt}
{\large 〔作者块由 submission.yaml 生成；匿名稿留空〕\par}
\end{center}
\draftnote{本文为实验前的中文稿件蓝图，与英文稿 \texttt{main.tex} 逐节同步（\texttt{paper\_assets/scripts/check\_zh\_sync.py} 机检）。
第 \ref{sec:problem}、\ref{sec:method} 节为正文；摘要、引言、相关工作、实验与结论为结构化草稿，论断标
\textsuperscript{[待裁决]}。摘要与贡献列表须在 \texttt{docs/claim-verdicts.md} 完成裁决后按实际结果重写，不得直接投稿。}
\section*{摘要}
\noindent
多重度柔性作业车间调度问题描述少数几类工件各有多个相同副本的车间：搜索空间随副本数增长，而同类副本可以互换。现有流体辅助的进化与强化学习方法把流体松弛当作初始化、状态特征或控制信号使用，没有说明离散调度的困难里哪一部分已被流体模型解释。本文把流体下界到可行 makespan 的差距精确分解为连续配额取整为整数配额时的取整损失与固定整数配额下的排序残差。维度不随副本数增长的配额染色体与守恒的配额迁移承载前者；只预测排序残差、输出不违反解析下界的商图排序器承载后者，并由学习器、流体启发式与随机三源并集的保底预选把每一次决策都交给精确解码裁决。\pending{在三个结构因子平衡交叉的 59 个固定设计算例、Brandimarte 公开基准与多重度外推档上，本方法在同墙钟预算下的 makespan gap 不高于 job-indexed 遗传算法、流体进化算法与两类学习型基线，真实解码次数减少 \OursDecodeReduction{}。}\pending{实证的取整差距随平均多重度减小，与理论预测一致。}

\vspace{6pt}
\noindent\textbf{关键词：}多重度柔性作业车间调度；流体松弛；整数配额；候选排序；进化搜索
\clearpage
\tableofcontents
\clearpage
% sections/zh/1-introduction.tex -- 漏斗五段式；与 sections/en/1-introduction.tex 逐段对应。
\section{引言}\label{sec:intro}

多重度柔性作业车间调度问题描述这样的车间：工件只有少数几种类型，每种类型有多个副本，副本共享同一条工艺路线、每道工序同一个候选机器集与同一组加工时长。电子装配、注塑与批量机加工车间都是这种形态：一个排产窗口内有几千道工序，但只有十来种产品。\pending{计划员每个窗口要出一版可执行的调度；用 job-indexed 遗传算法时，副本数翻倍后求解质量明显下滑。}难点在结构而非规模：同类副本可互换使 job-indexed 表示存在大量等价解，连续配额取整为整数带来的损失与固定配额下排序带来的损失来源不同。

\pending{多重度 FJSP 的流体松弛能给出有意义的分流比例与下界，流体辅助的进化与强化学习方法近年连续出现。}这些方法把流体解当作初始化、状态特征或控制信号，没有说明离散调度的困难里哪一部分已被流体模型解释；学习组件因此面对整个改进问题，在分布外失效时也没有非学习候选兜底。

本文从一个精确恒等式出发：流体下界到可行 makespan 的差距等于只依赖整数配额的取整损失与在该配额下由排序决定的排序残差之和。该恒等式给出了一种分工：进化搜索在维度由类型结构决定的配额染色体上进行，学习器只学习排序残差，精确解码是接受与否的唯一裁决者。

\draftnote{以下贡献列表为预期贡献。按反推式写作纪律，终稿的贡献列表只能由 \texttt{docs/claim-verdicts.md} 中裁决为“成立”的 claim 反推得到；“部分成立”的须收窄到证据边界内，“不成立”的必须删除。}

本文的贡献如下。
\begin{enumerate}
  \item 配额压缩表示及其流体-配额几何（第 \ref{sec:quota} 节）：基因型维度不随副本数增长，配额迁移保持各类型总量守恒，largest-remainder 配额的取整损失有与多重度无关的上界（命题 \ref{prop:rounding}）。
  \item 残差限定的商图排序器（第 \ref{sec:ranker} 节）：只预测排序残差，标签按流体下界归一化，输出不低于解析下界。
  \item 保底预选（第 \ref{sec:safeguard} 节）：学习器、流体启发式与随机三源并集在精确解码下裁决，局部不劣于单用启发式（命题 \ref{prop:safeguard}），并以可证伪的解码次数门槛评判。
  \item \pending{在 59 个固定设计算例、Brandimarte 公开基准与多重度外推档上的实验研究，含取整差距对其理论界的因子效应分析。}
\end{enumerate}

本文其余部分组织如下。第 \ref{sec:related} 节在流体辅助与学习辅助调度中定位本文；第 \ref{sec:problem} 节给出问题、流体松弛与差距恒等式；第 \ref{sec:method} 节给出方法及其结构性质；第 \ref{sec:experiments} 节报告实验；第 \ref{sec:conclusion} 节总结。

% sections/zh/2-related-work.tex -- 按缺口组织；与英文稿同一张定位表。
\section{相关工作}\label{sec:related}

\subsection{多重度 FJSP 的流体辅助进化搜索}\label{sec:related-fluid}
\pending{流体主从进化算法及其变体把多重度 FJSP 松弛为连续流体模型，用分流比例做初始化、调度规则或主种群的宏观指导\citeneed{流体主从进化算法（多重度 FJSP），2022，及其换模时间变体}。}这些工作证明流体比例携带有意义的分流结构，但流体解只是先验：搜索与学习组件仍面对完整的离散问题，表示也仍是 job-indexed。

\subsection{流体模型与强化学习的结合}\label{sec:related-rl}
\pending{以流体量作状态特征的多策略 PPO\citeneed{MPPPO，2024}、在线调节流体进化算法参数的 PPO\citeneed{DSLFMAE，2025}、把流体松弛送入双重 DQN 的动态多重度调度\citeneed{流体 + DDQN，2025}都把流体量当作特征或控制信号。}它们都没有分解差距的来源。\pending{一条并发工作研究副本数趋于无穷时流体解的精确可达性\citeneed{递归流体可达性，2026}}，与本文关注的有限、不均衡多重度区间互补。

\subsection{改进式学习与候选筛选学习}\label{sec:related-improve}
JSP 与 FJSP 的改进式深度强化学习让策略网络承担完整的改进决策并依赖信用分配\citeneed{L2S 类关键块改进策略，2024；FJSP 记忆增强改进启发式，2025}。监督式候选筛选，如 NeuroLKH 的候选边集与禁忌搜索的学习式移动筛选\citeneed{NeuroLKH；学习式禁忌移动筛选}，是本文排序器角色的方法论祖先，但未与基于松弛的分解结合。经典 FJSP 元启发式提供宿主组件：双层遗传编码\cite{pezzella2008genetic}、关键块禁忌搜索\cite{mastrolilli2000effective}，以及 Brandimarte\cite{brandimarte1993routing} 与 Hurink\cite{hurink1994tabu} 基准。

\begin{table}
\centering\small
\caption{最接近工作与本文方法的差异。}\label{tab:related}
\begin{tabularx}{\linewidth}{@{}>{\hsize=0.75\hsize}L>{\hsize=0.85\hsize}L>{\hsize=1.4\hsize}L@{}}
\toprule
\textbf{工作} & \textbf{范式} & \textbf{差异} \\
\midrule
流体主从进化算法 & 流体辅助进化 & 流体比例只作初始化与指导；表示为 job-indexed \\
DSLFMAE & 流体进化 + PPO 控制 & 学习器调节进化算法参数而非候选质量 \\
MPPPO & 流体特征 + DRL 构造 & 端到端构造，无搜索 \\
L2S 类改进式 DRL & 改进式 DRL & 策略承担完整的改进决策 \\
\textbf{本文} & \textbf{配额染色体上的进化 + 残差排序器} & \textbf{差距精确分解；学习器限定于残差；精确解码裁决} \\
\bottomrule
\end{tabularx}
\end{table}

表 \ref{tab:related} 给出定位。所列工作都未把流体差距分成来源不同的两部分，也未把学习器限定在解析模型无法解释的那一部分，这正是本文瞄准的缺口。

% sections/zh/3-problem.tex -- 问题描述、假设、符号、精确模型、流体松弛、差距恒等式；与英文稿逐段对应。
\section{问题描述与流体-配额几何}\label{sec:problem}

\subsection{问题描述}\label{sec:problem-desc}
车间含机器集合 $M$ 与工件类型集合 $\mathcal{R}$。类型 $r$ 有 $N_r \ge 1$ 个完全相同的副本，每个副本按固定工艺路线依次加工 $K_r$ 道工序。类型工序 $(r,k)$ 可在候选机器集 $M_{rk} \subseteq M$ 中任一机器上加工，在机器 $i$ 上的加工时长为 $p_{rki}$；同类型的副本共享候选机器集与加工时长，但每个副本须独立满足自身工艺路线的先后约束。决策为每个副本每道工序的机器分配与各机器上的加工顺序。全部数据在决策前已知，调度一次性给出，因此设定为静态确定性。$N_r = 1$ 对所有 $r$ 成立时，问题退化为标准柔性作业车间调度问题。

假设如下。A1：加工时长确定且已知（常规）；A2：机器同一时刻至多加工一道工序（常规）；A3：加工不可抢占（常规）；A4：所有副本在零时刻可用、无动态到达（常规，在局限中声明）；A5：忽略换模、运输与批量拆分（本文简化，在局限中声明）。

\subsection{符号}\label{sec:notation}
表 \ref{tab:notation} 列出符号，符号在全文只有一种含义，术语与译法汇总于附录表 \ref{tab:glossary}。副本工序记为 $o = (r,c,k)$，$c$ 为副本索引，$p_{oi} = p_{rki}$。类型的最短工艺路线长为 $\ell_r^{\min} = \sum_k \min_{i \in M_{rk}} p_{rki}$，$B$ 为足够大的常数。

\begin{table}[htbp]
\centering\small
\caption{符号说明}\label{tab:notation}
\begin{tabularx}{\linewidth}{@{}>{\hsize=0.72\hsize}L>{\hsize=1.28\hsize}L@{}}
\toprule
\textbf{符号} & \textbf{含义} \\
\midrule
$\mathcal{R},\,M$ & 工件类型集合、机器集合（$R = \lvert\mathcal{R}\rvert$，$m = \lvert M\rvert$） \\
$r,\,c,\,i$；$o,\,o'$ & 类型、副本、机器索引；副本工序 $o = (r,c,k)$ \\
$N_r,\,K_r$ & 类型 $r$ 的副本数；每副本工序数 \\
$M_{rk},\,p_{rki}$ & $(r,k)$ 的候选机器集与加工时长 \\
$\ell_r^{\min},\,B$ & 类型 $r$ 的最短工艺路线长；big-M 常数 \\
$\bar N,\,V_N,\,F,\,V_p$ & 平均多重度、多重度离散度、柔性度、工时异质度 \\
$x_{oi},\,s_o,\,y_{oo'}$ & 机器分配、开工时刻与先后次序变量 \\
$C_{\max}$ & 最大完工时间 \\
$q_{rki},\,L^F,\,\lambda_i$ & 流体变量、流体下界、机器负载对偶变量（只作特征） \\
$n_{rki},\,\kappa_{rk},\,\chi$ & 整数配额、优先键、个体 $\chi = (\mathbf{n}, \boldsymbol{\kappa})$ \\
$L^Q(\mathbf{n})$ & 整数配额下界 \\
$G_Q,\,G_S$；$g_Q,\,g_S$ & 取整损失、排序残差；二者按 $L^F$ 归一化的值 \\
$B_i,\,\bar B$ & 机器 $i$ 的取整界及其最大值 \\
$\mathcal{A},\,a,\,T_a$ & 候选池、候选动作、动作变换 \\
$\hat G_S(a)$ & 动作 $a$ 的预测排序残差 \\
$K_L,\,K_H,\,K_R$ & 学习器、流体启发式与随机三源各取的候选数 \\
\bottomrule
\end{tabularx}
\end{table}

\subsection{精确模型}\label{sec:model}
job-indexed 模型为
\begin{align}
\min\quad & C_{\max} \label{eq:obj}\\
\text{s.t.}\quad & \textstyle\sum_{i \in M_{rk}} x_{oi} = 1, && \forall\, o \label{eq:assign}\\
& s_{(r,c,k+1)} \ge s_o + \textstyle\sum_{i} p_{oi} x_{oi}, && \forall\, o:\, k < K_r \label{eq:prec}\\
& s_{o'} \ge s_o + p_{oi} - B(3{-}x_{oi}{-}x_{o'i}{-}y_{oo'}), && \forall\,i,\,o{\ne}o' \label{eq:disj}\\
& C_{\max} \ge s_o + \textstyle\sum_{i} p_{oi} x_{oi}, && \forall\, o:\, k = K_r \label{eq:cmax}\\
& x_{oi},\, y_{oo'} \in \{0,1\},\quad s_o \ge 0. \label{eq:dom}
\end{align}
式 \eqref{eq:assign} 把每道副本工序恰分配到一台候选机器，式 \eqref{eq:prec} 保证每个副本的工艺先后，式 \eqref{eq:disj} 是同机工序的 big-M 析取（$y_{oo'} = 1$ 表示 $o$ 在 $o'$ 之前），式 \eqref{eq:cmax} 定义 makespan，式 \eqref{eq:dom} 为变量域。该模型完整到可直接交给 CP-SAT\cite{ortools}，small 档的精确参考解即由此得到。

\subsection{配额型流体松弛}\label{sec:fluid}
松弛保留各类型的总量而放弃副本的离散性：
\begin{align}
L^F = \min\quad & L \label{eq:fluid-obj}\\
\text{s.t.}\quad & \textstyle\sum_{i \in M_{rk}} q_{rki} = N_r, && \forall\, r, k \label{eq:fluid-cons}\\
& \textstyle\sum_{r}\sum_{k:\, i \in M_{rk}} p_{rki}\, q_{rki} \le L, && \forall\, i \label{eq:fluid-load}\\
& L \ge \ell_r^{\min},\quad q_{rki} \ge 0, && \forall\, r, k, i. \label{eq:fluid-dom}
\end{align}
式 \eqref{eq:fluid-cons} 把类型工序 $(r,k)$ 的 $N_r$ 个副本按连续质量分流到候选机器，式 \eqref{eq:fluid-load} 使任一机器的流体负载不超过周期，式 \eqref{eq:fluid-dom} 使周期不小于任一类型的最短路线长。$L^F$ 是流体下界（命题 \ref{prop:fluidlb}），$q^*_{rki}/N_r$ 是 $(r,k)$ 的宏观分流比例。机器负载对偶变量 $\lambda_i$ 只作排序器特征：在流体最优解处它们近似均匀，几乎不含引导信息。

\subsection{整数配额与差距恒等式}\label{sec:gap}
整数配额 $\mathbf{n}$ 把 $(r,k)$ 的 $n_{rki}$ 个副本分配到机器 $i$，其中 $n_{rki}$ 为非负整数、$i \notin M_{rk}$ 时 $n_{rki} = 0$、$\sum_i n_{rki} = N_r$。它诱导配额下界
\begin{align}
L^Q(\mathbf{n}) &= \max\{ \max_i \textstyle\sum_{r,k:\, i \in M_{rk}} p_{rki} n_{rki},\, \max_r \ell_r^{\min} \}, \label{eq:LQ}\\
C_{\max}(\chi) - L^F &= \underbrace{L^Q(\mathbf{n}) - L^F}_{G_Q} + \underbrace{C_{\max} - L^Q(\mathbf{n})}_{G_S}. \label{eq:decomp}
\end{align}
恒等式 \eqref{eq:decomp} 是全文的锚点：任一可行个体的流体差距精确分解为只依赖配额的取整损失 $G_Q$ 与排序残差 $G_S$。归一化量 $g_Q = G_Q/L^F$、$g_S = G_S/L^F$ 用于跨算例比较与排序器标签。以下两条性质不依赖任何关于搜索的假设。

\begin{proposition}[流体下界]\label{prop:fluidlb}
任一满足式 \eqref{eq:assign}--\eqref{eq:dom} 的调度都有 $C_{\max} \ge L^F$。
\end{proposition}
\begin{proof}
令 $q_{rki} = \sum_c x_{(r,c,k),i}$。由式 \eqref{eq:assign} 得 $\sum_i q_{rki} = N_r$，即式 \eqref{eq:fluid-cons}。机器 $i$ 上的工序由式 \eqref{eq:disj} 互不重叠，总时长 $\sum_{r,k} p_{rki} q_{rki}$ 不超过 $C_{\max}$，即取 $L = C_{\max}$ 的式 \eqref{eq:fluid-load}。每个副本由式 \eqref{eq:prec} 串行完成 $K_r$ 道工序，故 $C_{\max} \ge \ell_r^{\min}$，式 \eqref{eq:fluid-dom} 成立。于是 $(q, C_{\max})$ 是式 \eqref{eq:fluid-obj}--\eqref{eq:fluid-dom} 的可行解，$L^F \le C_{\max}$。
\end{proof}

\begin{proposition}[非负分解]\label{prop:decomp}
对任一由配额个体 $\chi = (\mathbf{n}, \boldsymbol{\kappa})$ 解码得到的可行调度，式 \eqref{eq:decomp} 成立且 $G_Q(\chi) \ge 0$、$G_S(\chi) \ge 0$。
\end{proposition}
\begin{proof}
恒等式按定义成立。取 $q = \mathbf{n}$、$L = L^Q(\mathbf{n})$：式 \eqref{eq:fluid-cons} 由配额总量成立，式 \eqref{eq:fluid-load} 由式 \eqref{eq:LQ} 的第一项成立，式 \eqref{eq:fluid-dom} 由第二项成立，故 $L^Q(\mathbf{n}) \ge L^F$，$G_Q \ge 0$。解码后机器 $i$ 上的工序互不重叠，对所有 $i$ 有 $\sum_{r,k} p_{rki} n_{rki} \le C_{\max}$；每个副本串行完成全部工序，$C_{\max} \ge \ell_r^{\min}$；合起来 $C_{\max} \ge L^Q(\mathbf{n})$，$G_S \ge 0$。
\end{proof}

% sections/zh/4-method.tex -- 框架、三个机制、算法、性质、预测；与英文稿逐段对应。
\section{流体-残差进化搜索}\label{sec:method}

\subsection{总体框架}\label{sec:framework}
方法分四层（图 \ref{fig:framework}）。流体层对每个算例求解一次式 \eqref{eq:fluid-obj}--\eqref{eq:fluid-dom}，得到 $L^F$、$q^*$ 与 $\lambda^*$。配额层用 largest-remainder 取整把 $q^*$ 映射为整数配额 $\mathbf{n}^0$，构造初始种群，并由式 \eqref{eq:LQ}--\eqref{eq:decomp} 计算每个个体的 $g_Q$ 与 $g_S$。候选层对每个父代生成配额迁移与优先键动作的候选池 $\mathcal{A}$，取商图排序器、流体启发式与随机三源各自 top-$K$ 的并集。验证层对并集中的每个候选精确解码，按真实 makespan 选出子代并做精英选择。工作假设是：流体模型已经确定了负载与分流结构，整数配额只在有界的取整损失内偏离它（命题 \ref{prop:rounding}），剩余的自由度在排序残差；限定在该残差上的学习器面对更小、更可迁移的任务，而精确解码与保底并集把学习器失效的代价封顶（命题 \ref{prop:safeguard}）。

\begin{figure}[htbp]
\figplace{总体框架图：左侧流体层（配额型流体 LP，输出 $L^F$、$q^*$）与配额层（largest-remainder 取整 → 配额染色体 $(\mathbf{n}, \boldsymbol{\kappa})$；式 \eqref{eq:decomp} 的 $g_Q$ / $g_S$ 仪表）；中间候选层（配额迁移与优先键动作的候选池 → 排序器 / 流体启发式 / 随机三源 top-$K$ 并集，用颜色标出 C1 配额几何、C2 排序器、C3 保底并集）；右侧验证层（对称感知解码器 → 真实 $C_{\max}$ → $\mu+\lambda$ 精英选择）；回流箭头表示自举排序器训练标签的搜索轨迹}{无（示意图）}{opt-paper-scifig 路线 B（TikZ），双栏宽 17.8\,cm}{6cm}
\caption{流体-残差进化搜索：流体层、配额层、候选层与验证层}\label{fig:framework}
\end{figure}

\subsection{配额压缩表示与流体-配额几何}\label{sec:quota}
个体为 $\chi = (\mathbf{n}, \boldsymbol{\kappa})$：类型工序到机器的整数配额张量 $\mathbf{n}$，以及每个类型工序一个优先键 $\kappa_{rk}$。不为任何副本保留基因；副本在解码时展开，同类副本之间的对称性由此消除。

\paragraph{对称感知事件解码器} 解码器按时间推进。机器 $i$ 空闲时，它考察满足 $n_{rki} > 0$ 且至少有一个副本已完成第 $k-1$ 道工序的类型工序 $(r,k)$，按优先键选出最优的一类，再从该类型满足条件的副本中选最早可释放者。开工时刻取该副本前序完工与机器就绪的较大者，配额 $n_{rki}$ 减一。

\begin{theorem}[解码器正确性]\label{thm:decoder}
设 $\mathbf{n}$ 为非负整数张量，$i \notin M_{rk}$ 时 $n_{rki} = 0$ 且 $\sum_i n_{rki} = N_r$，$\boldsymbol{\kappa}$ 任意。则 (i) 事件解码器恰在 $\sum_r N_r K_r$ 次放置后终止，输出满足式 \eqref{eq:assign}--\eqref{eq:dom} 的调度；(ii) 当 $N_r = 1$ 对所有 $r$ 成立时，$\mathbf{n}$ 是机器分配向量，$\boldsymbol{\kappa}$ 诱导工序排列，解码器即标准 FJSP 的主动解码器。
\end{theorem}
\begin{proof}
(i) 只在“$(r,k)$ 在机器 $i$ 上尚有配额且某副本的前序已完成”时放置一道副本工序，每次放置使一个配额减一、一个副本前进一道工序。配额总和为 $\sum_r N_r K_r$，故解码器恰在这么多次放置后停止，且每道副本工序恰得到一台候选机器，即式 \eqref{eq:assign}。开工时刻不早于前序完工，即式 \eqref{eq:prec}；同机工序按放置顺序追加且互不重叠，取 $y$ 为该顺序即满足式 \eqref{eq:disj}；式 \eqref{eq:cmax} 与 \eqref{eq:dom} 按构造成立。非终止状态总存在可放置项：每个未完成副本的首道未完成工序的前序已完成，而该类型工序的剩余配额等于尚未执行它的副本数，故必有某台机器的配额为正。(ii) $N_r = 1$ 时 $n_{rki} \in \{0,1\}$ 且每个 $(r,k)$ 恰一个单位，即机器分配向量；副本的选择平凡，解码器按优先键在就绪工序中选择并以最早可行时刻开工，即主动解码。
\end{proof}

\paragraph{流体-配额几何} 流体解通过两条路径进入搜索。初始化时 largest-remainder 取整把 $q^*$ 映射为 $\mathbf{n}^0$；搜索中配额迁移把 $n_{rk\cdot}$ 的 $b$ 个单位从配额超过 $q^*_{rki}$ 的机器移到配额不足的机器，平局随机打破，于是 $\sum_i n_{rki} = N_r$ 守恒，$L^Q$ 由式 \eqref{eq:LQ} 增量更新。与流体辅助进化算法不同，流体解定义的是表示的坐标系与动作的方向而不只是初始化；与 job-indexed 编码不同，基因型维度不含 $N_r$。

\begin{proposition}[取整界]\label{prop:rounding}
设 $q^*$ 为式 \eqref{eq:fluid-obj}--\eqref{eq:fluid-dom} 的最优解，$\mathbf{n}^0$ 为对每个 $q^*_{rk\cdot}$ 按 largest-remainder 法取整的配额，$B_i = \sum_{(r,k):\, i \in M_{rk}} p_{rki}$，$\bar B = \max_i B_i$。则 $0 \le L^Q(\mathbf{n}^0) - L^F \le \bar B$，且 $\bar B$ 不依赖任何 $N_r$。
\end{proposition}
\begin{proof}
下界即命题 \ref{prop:decomp}。largest-remainder 取整保持 $\sum_i n^0_{rki} = N_r$ 且 $\lvert n^0_{rki} - q^*_{rki}\rvert < 1$。对机器 $i$，由式 \eqref{eq:fluid-load}，$\sum_{r,k} p_{rki} n^0_{rki} < \sum_{r,k} p_{rki} q^*_{rki} + \sum_{(r,k):\, i \in M_{rk}} p_{rki} \le L^F + B_i$；又由式 \eqref{eq:fluid-dom}，$\max_r \ell_r^{\min} \le L^F$。代入式 \eqref{eq:LQ} 得 $L^Q(\mathbf{n}^0) \le L^F + \bar B$。每个 $B_i$ 只对类型工序求和。
\end{proof}

\begin{proposition}[商图维度]\label{prop:dims}
配额染色体有 $\sum_r K_r \lvert M_{rk}\rvert + \sum_r K_r$ 个分量；商图有 $\sum_r K_r + \lvert M\rvert$ 个节点、$\sum_r (K_r - 1) + \sum_r K_r \lvert M_{rk}\rvert$ 条边；三者都不依赖 $N_r$。job-indexed 双层编码有 $2\sum_r N_r K_r$ 个分量。
\end{proposition}
\begin{proof}
按定义计数：每个 $i \in M_{rk}$ 的 $(r,k,i)$ 一个整数，每个 $(r,k)$ 一个优先键；商图节点为类型工序与机器，边为同类型相邻工序对与类型工序-候选机器对。job-indexed 编码为每道副本工序各保留一个机器基因与一个序列位置。
\end{proof}

\subsection{残差限定的商图排序器}\label{sec:ranker}
排序器在命题 \ref{prop:dims} 的商图上工作。类型工序节点携带 $N_r$、$K_r$、$q^*_{rki}/N_r$、$n_{rki}/N_r$、配额超额量、加工时长统计、当前关键性与优先键统计；机器节点携带 $\lambda^*_i$、流体与配额负载占比，以及解码器中观测到的空闲占比与关键负载占比；全局特征为 $g_Q$、$g_S$、$L^Q/L^F$、$\bar N$ 与 $V_N$。对候选动作 $a$，网络由其涉及节点的嵌入与动作类型形成候选嵌入，输出 $\hat z_a$ 并经
\begin{equation}
\hat G_S(a) = L^F \operatorname{softplus}(\hat z_a) \label{eq:softplus}
\end{equation}
映射，使预测的子代 makespan $L^Q(\mathbf{n}_a) + \hat G_S(a) \ge L^F$ 不低于解析下界。标签 $y_a = G_S(T_a\chi)/L^F$ 由无排序器的宿主 FRES-lite 的搜索轨迹上精确解码得到，损失为同一父代内的 pairwise logistic 排序损失。排序器既不选择也不接受动作，不需要强化学习的信用分配；与通用 move ranker 不同，它的标签、特征与输出都限定在式 \eqref{eq:decomp} 的第二项上，解析已知的 $L^Q(\mathbf{n}_a)$ 不再学习。

\subsection{保底预选}\label{sec:safeguard}
每个父代的候选池 $\mathcal{A}$ 中配额迁移与优先键动作（关键块交换、关键工序插入与键扰动）各约一半。只对一个小子集精确解码：
\begin{equation}
\mathcal{A}_{\mathrm{eval}} = \operatorname{Top}_{K_L}\bigl(\hat G_S\bigr) \cup \operatorname{Top}_{K_H}(\text{流体启发式}) \cup \operatorname{Rand}_{K_R}, \label{eq:union}
\end{equation}
其中流体启发式按 $\sum_i \lvert n_{rki} - q^*_{rki}\rvert$ 的减少量给配额迁移打分，按所涉工序是否在关键块中给键动作打分。子代取解码后 makespan 最小的候选。

\begin{proposition}[保底局部不劣]\label{prop:safeguard}
设 $\mathcal{A}_H \subseteq \mathcal{A}$ 为启发式候选集，$\mathcal{A}_{\mathrm{eval}} = \mathcal{A}_L \cup \mathcal{A}_H \cup \mathcal{A}_R$，全部精确解码。则 $\min_{a \in \mathcal{A}_{\mathrm{eval}}} C_{\max}(T_a\chi) \le \min_{a \in \mathcal{A}_H} C_{\max}(T_a\chi)$。
\end{proposition}
\begin{proof}
$\mathcal{A}_H \subseteq \mathcal{A}_{\mathrm{eval}}$，最小值在集合包含下单调不增。该陈述不涉及排序器的推断开销，墙钟上的优势是实验主张（第 \ref{sec:ablation} 节）。
\end{proof}

\subsection{算法、训练与一致性}\label{sec:algorithm}
算法 \ref{alg:fres} 为宿主搜索；算法 \ref{alg:train} 在按算例参数表随机构造的算例上离线训练排序器并冻结，因此每个评测算例都是零样本求解。

\begin{algorithm}[htbp]
\caption{FRES：流体-残差进化搜索}\label{alg:fres}
\begin{algorithmic}[1]
\Require 算例；种群 $\mu$、子代 $\lambda$；候选池大小 $B$；$K_L, K_H, K_R$；配额步长 $b$；排序器 $f$（FRES-lite 无）；墙钟预算 $T$
\Ensure 最优解码调度
\State 求解式 \eqref{eq:fluid-obj}--\eqref{eq:fluid-dom}：$L^F, q^*, \lambda^*$；$\mathbf{n}^0 \gets$ $q^*$ 的 largest-remainder 取整
\State $P \gets \mu$ 个个体（$\mathbf{n}^0$ 加随机迁移，$\boldsymbol{\kappa}$ 随机）；逐个解码，记录 $C_{\max}, G_Q, G_S$
\While{已用时间 $< T$}
  \For{$t = 1$ \textbf{to} $\lambda$}
    \State $\chi \gets$ 从 $P$ 锦标赛选出的父代
    \State $\mathcal{A} \gets$ $B$ 个候选的池：配额迁移（超额机器 $\to$ 欠额机器，步长 $b$）与优先键动作
    \State $\mathcal{A}_H \gets$ 按流体启发式取 $\mathcal{A}$ 的 $\operatorname{Top}_{K_H}$；若给定 $f$，$\mathcal{A}_L \gets$ 按 $L^Q(\mathbf{n}_a) + f(a)$ 取 $\operatorname{Top}_{K_L}$，否则 $\emptyset$；$\mathcal{A}_R \gets$ 其余候选中的 $\operatorname{Rand}_{K_R}$
    \ForAll{$a \in \mathcal{A}_L \cup \mathcal{A}_H \cup \mathcal{A}_R$}
      \State 精确解码 $T_a\chi$；记录 $C_{\max}$ 与 $G_S$（收集标签时记录 $(\chi, a, G_S/L^F)$）
    \EndFor
    \State 把 $\arg\min_a C_{\max}(T_a\chi)$ 加入子代集合
  \EndFor
  \State $P \gets$ 按 $C_{\max}$ 的 $\mu+\lambda$ 精英选择，附多样性救援比例
\EndWhile
\State \Return $P$ 中的最优个体
\end{algorithmic}
\end{algorithm}

\begin{algorithm}[htbp]
\caption{残差排序器的离线训练}\label{alg:train}
\begin{algorithmic}[1]
\Require 算例参数表；epoch 数 $E$；每 epoch 算例数 $n_e$；学习率；验证集
\Ensure 冻结的排序器 $f$
\For{$e = 1$ \textbf{to} $E$}
  \State 随机构造 $n_e$ 个算例；运行不带 $f$ 的算法 \ref{alg:fres}（FRES-lite），记录三元组 $(\chi, a, G_S(T_a\chi)/L^F)$
  \State 构建商图；用同一父代内候选对的 pairwise logistic 排序损失更新 $f$
  \State 用当前 $f$ 的完整 FRES 在验证集上评测；保留最优检查点；平台期早停
\EndFor
\State \Return 最优检查点
\end{algorithmic}
\end{algorithm}

\begin{proposition}[学习目标与模型目标一致]\label{prop:consistency}
设 $\mathbf{n}_a$ 为子代 $T_a\chi$ 的配额（迁移动作确定性地改写配额，键动作保持配额不变，故 $\mathbf{n}_a$ 无需解码即可得到）。若 $\hat G_S(a) = G_S(T_a\chi)$ 对所有 $a \in \mathcal{A}$ 成立，则候选按 $L^Q(\mathbf{n}_a) + \hat G_S(a)$ 的排序与按真实 $C_{\max}(T_a\chi)$ 的排序一致；以 $G_S(T_a\chi)/L^F$ 为标签的 pairwise 排序损失在同一父代内取零当且仅当预测保持 $G_S$ 的序。
\end{proposition}
\begin{proof}
对 $T_a\chi$ 应用命题 \ref{prop:decomp}，$C_{\max}(T_a\chi) = L^F + G_Q(T_a\chi) + G_S(T_a\chi) = L^Q(\mathbf{n}_a) + G_S(T_a\chi)$。预测精确时两种排序键逐项相等。同一父代内 $L^F$ 为正常数，除以它不改变序，而 pairwise 损失只惩罚序被反转的候选对。
\end{proof}

\paragraph{复杂度} 流体 LP 有 $\sum_r K_r \lvert M_{rk}\rvert + 1$ 个变量，只求解一次。每机器一个堆时一次解码的代价为 $O(\sum_r N_r K_r \log m)$；每代花费 $\lambda (K_L + K_H + K_R)$ 次解码，外加排序器对 $\lambda B$ 个候选的一次批量前向，其图规模由命题 \ref{prop:dims} 给出。job-indexed 宿主把同样的解码次数花在大小为 $2\sum_r N_r K_r$ 的基因型上。

\subsection{待实验检验的预测}\label{sec:predictions}
性能论断不写成定理，而是写成猜想并附证伪条件。

\begin{conjecture}[取整差距的符号]\label{conj:rounding}
在第 \ref{sec:settings} 节的算例分布下，实证的相对取整差距 $G_Q(\mathbf{n}^0)/L^F$ 随 $\bar N$ 增大而减小、随 $V_p$ 增大而增大。
\end{conjecture}
命题 \ref{prop:rounding} 的界为 $\bar B/L^F$；$L^F$ 随 $\bar N$ 近似线性增长而 $\bar B$ 不变，$V_p$ 增大时 $B_i$ 中的大项与流体分流偏爱的小项分离。曲线在 $\bar N$ 上不单调或在 $V_p$ 上符号反转即推翻本猜想。

\begin{conjecture}[增益的符号]\label{conj:gain}
在主实验档，FRES 相对最强基线的 gap 增益随 $\bar N$ 与 $F$ 增大、随 $V_N$ 增大而减小。
\end{conjecture}
$\bar N$ 增大时 job-indexed 表示的对称冗余增大而配额表示不变（命题 \ref{prop:dims}）；$F$ 增大时配额自由度与流体分流的信息量增大；$V_N$ 增大时稀有类型相对其流体质量的取整损失更大。任一符号反转即推翻本猜想，并在第 \ref{sec:factor} 节如实报告。

% sections/zh/5-experiments.tex -- 按规格 6C 的结构化骨架（5.1–5.7 + 算例因子效应）；数值全部为占位符。
\section{实验}\label{sec:experiments}

\draftnote{本节全部数值为占位符，将由 \texttt{paper\_assets/scripts/make\_tables.py} 从 \texttt{result/*.csv} 按《实验设计规格》6B 契约填充；全部图形为占位框，由 opt-paper-scifig 生成后替换。方向性表述在裁决完成前均为预期。}

\subsection{实验设置}\label{sec:settings}
评测集是由算例参数表生成的固定设计集：平均多重度 $\bar N$ 为规模因子，多重度离散度 $V_N$、柔性度 $F$ 与工时异质度 $V_p$ 为三个结构因子，类型数 $R = 8$、机器数 $m = 10$（表 \ref{tab:instances}）。主实验档为 $\bar N = 10$ 下的 $3\times3\times3$ 全因子；small 档由 CP-SAT\cite{ortools} 在 1800 s 内给出精确参考解；large 档把多重度外推到训练区间之外；Brandimarte 算例\cite{brandimarte1993routing} 为 $N_r = 1$ 的特例。算例以文件形式固定并随论文发布，复现依靠这些文件与 \NumRuns{} 次独立运行而非随机种子。表 \ref{tab:hyper} 列出超参数，表 \ref{tab:baselines} 列出基线：每族一个强代表，取推荐配置与同等墙钟预算。指标为 makespan、相对精确解或参考解的最优性 gap、用于跨档比较的流体 gap $(C_{\max} - L^F)/L^F$、可行率、1 / 10 / 60 s 检查点下的质量-时间曲线面积，以及精确解码次数与达到目标的时间。显著性用逐算例 run 均值的配对 Wilcoxon 检验（$p < 0.05$）评估，方法数不少于四时另给 Friedman–Nemenyi 临界差图。

\begin{table}[htbp]
\centering\small
\caption{评测算例设置}\label{tab:instances}
% generated by paper_assets/scripts/make_tables.py -- do not edit by hand

\begin{tabularx}{\linewidth}{@{}l>{\hsize=1.000\hsize\raggedright\arraybackslash}X>{\hsize=0.900\hsize\raggedright\arraybackslash}Xl>{\hsize=1.100\hsize\raggedright\arraybackslash}X@{}}
\toprule
\textbf{\bi{档位}{Tier}} & \textbf{\bi{规模因子 $\bar N$}{Scale factor $\bar N$}} & \textbf{\bi{结构因子}{Structural factors}} & \textbf{\bi{算例数}{Instances}} & \textbf{\bi{用途}{Purpose}} \\
\midrule
Small & 2, 4 & \bi{各取低 / 高两端}{low / high ends} & 16 & \bi{精确参考解校准真实 gap}{exact reference solutions calibrate the true gap} \\
Main & 10 & \bi{3 $\times$ 3 $\times$ 3 全因子}{$3\times3\times3$ full factorial} & 27 & \bi{主对比 + 因子效应分析}{main comparison and factor-effect analysis} \\
Large & 25, 50 & \bi{各取低 / 高两端}{low / high ends} & 16 & \bi{多重度外推（商图表示直接求解）}{multiplicity extrapolation (quotient representation solved directly)} \\
Benchmark & \bi{Brandimarte mk01–mk10（$N_r = 1$ 退化为标准 FJSP）}{Brandimarte mk01--mk10 ($N_r = 1$, standard FJSP)} & \bi{外部给定}{given} & 10 & \bi{公开基准（补充，不替代设计集）}{public benchmark (supplementary to the design set)} \\
\bottomrule
\end{tabularx}

\end{table}

\begin{table}[htbp]
\centering\small
\caption{超参数设置（草案值；终值以确认后的训练参数表为准）}\label{tab:hyper}
% generated by paper_assets/scripts/make_tables.py -- do not edit by hand

\begin{tabularx}{\linewidth}{@{}l>{\hsize=0.850\hsize\raggedright\arraybackslash}X>{\hsize=1.150\hsize\raggedright\arraybackslash}X@{}}
\toprule
\textbf{\bi{类别}{Group}} & \textbf{\bi{参数}{Parameter}} & \textbf{\bi{取值}{Value}} \\
\midrule
\bi{训练}{Training} & \texttt{train.paradigm} & \bi{generalist（排序器离线训练、冻结后零样本）}{generalist (ranker trained offline, frozen, zero-shot)} \\
\bi{训练}{Training} & \texttt{instances\_\allowbreak{}per\_\allowbreak{}epoch} & \bi{256 个随机算例，FRES-lite 自举 32\,768 个候选-标签对}{256 random instances; 32\,768 candidate--label pairs bootstrapped by FRES-lite} \\
\bi{训练}{Training} & \texttt{epochs} /\allowbreak{} \texttt{plateau\_\allowbreak{}patience} & 60 / 10 \\
\bi{训练}{Training} & \texttt{wall\_\allowbreak{}budget\_\allowbreak{}s} /\allowbreak{} \texttt{eval\_\allowbreak{}every} /\allowbreak{} \texttt{env\_\allowbreak{}batch} & \bi{待实测}{to be measured} \\
\bi{算法}{Optimizer} & \texttt{lr} & \bi{1e-3，余弦退火}{1e-3, cosine annealing} \\
\bi{算法}{Optimizer} & \texttt{batch\_\allowbreak{}size} /\allowbreak{} \texttt{loss} & \bi{512 / pairwise ranking（logistic）}{512 / pairwise ranking (logistic)} \\
\bi{网络}{Network} & \texttt{embed\_\allowbreak{}dim} /\allowbreak{} \texttt{n\_\allowbreak{}layers} & \bi{128 / 3（异构消息传递）}{128 / 3 (heterogeneous message passing)} \\
\bi{宿主}{Host EA} & \texttt{population} /\allowbreak{} \texttt{offspring} & 100 / 100 \\
\bi{宿主}{Host EA} & $B$ / $K_L$ / $K_H$ / $K_R$ & 96 / 10 / 5 / 3 \\
\bi{宿主}{Host EA} & \bi{配额迁移步长 $b$ / 优先键扰动幅度}{quota step $b$ / priority-key perturbation} & 1 / 0.1 \\
\bottomrule
\end{tabularx}

\end{table}

\begin{table}[htbp]
\centering\small
\caption{对比方法及配置}\label{tab:baselines}
% generated by paper_assets/scripts/make_tables.py -- do not edit by hand

\begin{tabularx}{\linewidth}{@{}>{\hsize=0.700\hsize\raggedright\arraybackslash}Xl>{\hsize=1.300\hsize\raggedright\arraybackslash}X@{}}
\toprule
\textbf{\bi{方法}{Method}} & \textbf{\bi{类别}{Class}} & \textbf{\bi{实现与配置}{Implementation and configuration}} \\
\midrule
\bi{HGA（job-indexed）}{HGA (job-indexed)} & \bi{元启发式}{metaheuristic} & \bi{文献推荐参数（双层编码遗传算法）；种群 100、POX 交叉 0.8、变异 0.15、精英 10\%；同墙钟；每算例 3 次重启}{recommended settings of the two-layer encoded GA; population 100, POX crossover 0.8, mutation 0.15, elitism 10\%; same wall clock; 3 restarts per instance} \\
\bi{关键块禁忌搜索（TS）}{Critical-block tabu search (TS)} & \bi{元启发式}{metaheuristic} & \bi{Mastrolilli–Gambardella 型，自实现并在 mk 基准核对；禁忌长度 12、关键块邻域 + 关键工序改派；同墙钟；3 次重启}{Mastrolilli--Gambardella type, reimplemented and checked on the mk benchmark; tenure 12, critical-block neighbourhood with reassignment; same wall clock; 3 restarts} \\
\bi{流体进化算法（FMAE-s）}{Fluid-assisted EA (FMAE-s)} & \bi{元启发式}{metaheuristic} & \bi{流体主从进化算法的无换模静态版（复现）；原文推荐参数；流体比例只作初始化与规则；同墙钟；3 次重启}{static, setup-free version of the fluid master--slave EA (reimplementation); published settings; fluid ratios used only for initialization and rules; same wall clock; 3 restarts} \\
DSLFMAE & \bi{学习型}{learning-based} & \bi{PPO 辅助流体进化算法（复现或机制等价复现，标注 reimplementation）；原文推荐配置；与本方法同等训练预算，3 run}{PPO-assisted fluid EA (reimplementation, marked as such); published configuration; same training budget as ours, 3 runs} \\
HGNN-DRL & \bi{学习型}{learning-based} & \bi{官方代码（异构图神经网络 + PPO 构造式调度）；推荐配置；贪心 + 采样两口径；同等训练预算，3 run}{official code (heterogeneous GNN with PPO, constructive); recommended configuration; greedy and sampling decoding; same training budget, 3 runs} \\
CP-SAT & \bi{参考解}{reference} & \bi{OR-Tools；1800 s（small）/ 与元启发式同墙钟（main，anytime 参考）}{OR-Tools; 1800 s (small) / same wall clock as the metaheuristics (main, anytime reference)} \\
\bottomrule
\end{tabularx}

\end{table}

\subsection{与基线方法的对比}\label{sec:comparison}
本小节回答 FRES 是否优于元启发式与学习型基线、差异是否显著。表 \ref{tab:main} 报告 small 与 main 档在 10 s 与 60 s 检查点的结果；图 \ref{fig:main} 给出 gap 分布与临界差图。\pending{仅 FRES-lite 就达到不高于 HGA 与 FMAE-s 的 gap，完整方法在主实验档相对最强基线改善 \ImpvOverBest{}。}

\begin{table}[htbp]
\centering\small
\caption{主要结果：\NumRuns{} 次运行的均值 $\pm$ 标准差；最优加粗；$\dagger$ 表示显著优于次优（Wilcoxon，$p < 0.05$）。S = small 档（相对精确解的 gap），M = main 档（相对参考解的 gap）}\label{tab:main}
% generated by paper_assets/scripts/make_tables.py -- do not edit by hand

\begin{tabularx}{\linewidth}{@{}>{\hsize=1.500\hsize\raggedright\arraybackslash}X>{\hsize=0.800\hsize\raggedright\arraybackslash}X>{\hsize=0.800\hsize\raggedright\arraybackslash}X>{\hsize=1.000\hsize\raggedright\arraybackslash}X>{\hsize=1.000\hsize\raggedright\arraybackslash}X>{\hsize=0.800\hsize\raggedright\arraybackslash}X>{\hsize=1.100\hsize\raggedright\arraybackslash}X@{}}
\toprule
\textbf{\bi{方法}{Method}} & \textbf{\bi{S-gap（\%）}{S gap (\%)}} & \textbf{\bi{S-可行率}{S feas.}} & \textbf{\bi{M-gap 10 s（\%）}{M gap 10 s (\%)}} & \textbf{\bi{M-gap 60 s（\%）}{M gap 60 s (\%)}} & \textbf{\bi{M-可行率}{M feas.}} & \textbf{\bi{解码次数（$\times10^{3}$）}{Decodings ($\times10^{3}$)}} \\
\midrule
\bi{HGA（job-indexed）}{HGA (job-indexed)} & \phpm & \ph & \phpm & \phpm & \ph & \ph \\
\bi{关键块禁忌搜索（TS）}{Critical-block tabu search (TS)} & \phpm & \ph & \phpm & \phpm & \ph & \ph \\
\bi{流体进化算法（FMAE-s）}{Fluid-assisted EA (FMAE-s)} & \phpm & \ph & \phpm & \phpm & \ph & \ph \\
\bi{DSLFMAE}{DSLFMAE} & \phpm & \ph & \phpm & \phpm & \ph & \ph \\
\bi{HGNN-DRL}{HGNN-DRL} & \phpm & \ph & \phpm & \phpm & \ph & \ph \\
\bi{CP-SAT（同墙钟）}{CP-SAT (same wall clock)} & \phpm & \ph & \phpm & \phpm & \ph & \ph \\
\bi{\textbf{FRES}}{\textbf{FRES}} & \phb & \phb & \phb & \phb & \phb & \phb \\
\bottomrule
\end{tabularx}

\end{table}

\begin{figure}[htbp]
\figplace{主实验档 60 s 时各方法的 gap 箱线图，每档一个面板，带 Wilcoxon 显著性标注；内嵌临界差图}{\texttt{result/*/eval\_results.csv}（gap、method、tier）；$p$ 值与秩取自 \texttt{result/stats\_summary.csv}}{opt-paper-scifig 路线 A（matplotlib），栏宽 88\,mm}{52mm}
\caption{主实验档的 gap 分布与临界差图}\label{fig:main}
\end{figure}

\subsection{消融与排序研究}\label{sec:ablation}
本小节回答增益是否来自三个机制、排序器是否减少了无效解码。变体 A1--A5 逐一去掉流体配额、配额压缩、排序器、残差限定与保底（表 \ref{tab:ablation}）。图 \ref{fig:ablation} 给出各变体的收敛过程，图 \ref{fig:recall} 给出同候选预算下排序器相对流体启发式与通用排序器的 top-10 recall，图 \ref{fig:safeguard} 给出被接受子代来自并集各源的占比。\pending{完整方法比 FRES-lite 少解码 \OursDecodeReduction{} 个候选而质量不降。}

\begin{table}[htbp]
\centering\small
\caption{主实验档的消融与机制指标}\label{tab:ablation}
% generated by paper_assets/scripts/make_tables.py -- do not edit by hand

\begin{tabularx}{\linewidth}{@{}>{\hsize=1.000\hsize\raggedright\arraybackslash}Xllll@{}}
\toprule
\textbf{\bi{变体}{Variant}} & \textbf{\bi{gap（\%）}{gap (\%)}} & \textbf{\bi{可行率}{feas.}} & \textbf{top-10 recall} & \textbf{\bi{解码次数（$\times10^{3}$）}{Decodings ($\times10^{3}$)}} \\
\midrule
\bi{完整方法（FRES）}{Full method (FRES)} & \phb & \phb & \phb & \phb \\
\bi{A1 无流体配额}{A1 no fluid quota} & \phpm & \ph & \ph & \ph \\
\bi{A2 无配额压缩}{A2 no quota compression} & \phpm & \ph & \ph & \ph \\
\bi{A3 无排序器（FRES-lite）}{A3 no ranker (FRES-lite)} & \phpm & \ph & \ph & \ph \\
\bi{A4 通用排序器}{A4 generic ranker} & \phpm & \ph & \ph & \ph \\
\bi{A5 无保底}{A5 no safeguard} & \phpm & \ph & \ph & \ph \\
\bottomrule
\end{tabularx}

\end{table}

\begin{figure}[htbp]
\figplace{FRES 与 A1--A5 的在位最优 gap 随墙钟变化，均值线加 $\pm$ 标准差阴影带，时间轴取对数}{\texttt{result/*/curve.csv} 跨 run 聚合}{opt-paper-scifig 路线 A，栏宽 88\,mm}{50mm}
\caption{消融变体的收敛过程}\label{fig:ablation}
\end{figure}

\begin{figure}[htbp]
\figplace{真实改善动作的 top-10 recall 随解码候选数变化，三条曲线（残差排序器、流体启发式、通用排序器），同一候选池}{\texttt{result/*/log.csv} 的 ranking 指标与 \texttt{operator\_trace.csv}}{opt-paper-scifig 路线 A，栏宽 88\,mm}{50mm}
\caption{同候选预算下的排序质量}\label{fig:recall}
\end{figure}

\begin{figure}[htbp]
\figplace{被接受子代的来源（学习器 / 流体启发式 / 随机）随代数变化的堆叠柱，每个 $V_N$ 水平一个面板}{\texttt{operator\_trace.csv}（best\_from、generation）join \texttt{index.csv}}{opt-paper-scifig 路线 A，栏宽 88\,mm}{48mm}
\caption{保底并集下被接受子代的来源}\label{fig:safeguard}
\end{figure}

\subsection{参数敏感性}\label{sec:sensitivity}
本小节回答预选参数是否需要调参。图 \ref{fig:sens} 在验证档上扫描候选池大小 $B$ 与学习器份额 $K_L$。\pending{性能在扫描网格上平坦。}

\begin{figure}[htbp]
\figplace{主实验档 gap 在 $B \in \{48, 96, 192\}$ 与 $K_L \in \{8, 16, 24\}$ 上的热力图，格内标均值，选定设置加框}{\texttt{result/sensitivity.csv}}{opt-paper-scifig 路线 A，栏宽 88\,mm}{48mm}
\caption{对候选池大小与学习器份额的敏感性}\label{fig:sens}
\end{figure}

\subsection{算例因子效应}\label{sec:factor}
本小节回答方法在什么问题特性下更强或更弱、理论界是否得到支持。它与上一小节的参数敏感性不同：这里的因子属于问题而非算法。图 \ref{fig:rounding} 在 $\bar N$ 与 $V_p$ 上比较实证取整差距 $G_Q/L^F$ 与命题 \ref{prop:rounding} 的界 $\bar B/L^F$（猜想 \ref{conj:rounding}）；表 \ref{tab:factor} 与图 \ref{fig:factor} 给出相对最强基线增益的主效应与双因子交互（猜想 \ref{conj:gain}）。误差棒为运行间标准差，不是单元内重复。

\begin{table}[htbp]
\centering\small
\caption{主实验档的算例因子主效应（27 个参数各异的算例，每水平 9 个；误差为 \NumRuns{} 次运行的运行间标准差）}\label{tab:factor}
% generated by paper_assets/scripts/make_tables.py -- do not edit by hand

\begin{tabularx}{\linewidth}{@{}>{\hsize=0.900\hsize\raggedright\arraybackslash}Xll>{\hsize=1.100\hsize\raggedright\arraybackslash}Xl@{}}
\toprule
\textbf{\bi{因子}{Factor}} & \textbf{\bi{水平}{Level}} & \textbf{\bi{gap（\%）}{gap (\%)}} & \textbf{\bi{相对最强基线的增益（pp）}{Gain over the strongest baseline (pp)}} & \textbf{\bi{解码次数（$\times10^{3}$）}{Decodings ($\times10^{3}$)}} \\
\midrule
\bi{多重度离散度 $V_N$}{Multiplicity dispersion $V_N$} & 0.0 & \phpm & \ph & \ph \\
\bi{多重度离散度 $V_N$}{Multiplicity dispersion $V_N$} & 0.5 & \phpm & \ph & \ph \\
\bi{多重度离散度 $V_N$}{Multiplicity dispersion $V_N$} & 1.0 & \phpm & \ph & \ph \\
\bi{柔性度 $F$}{Flexibility $F$} & 0.3 & \phpm & \ph & \ph \\
\bi{柔性度 $F$}{Flexibility $F$} & 0.55 & \phpm & \ph & \ph \\
\bi{柔性度 $F$}{Flexibility $F$} & 0.8 & \phpm & \ph & \ph \\
\bi{工时异质度 $V_p$}{Processing-time heterogeneity $V_p$} & 0.2 & \phpm & \ph & \ph \\
\bi{工时异质度 $V_p$}{Processing-time heterogeneity $V_p$} & 0.5 & \phpm & \ph & \ph \\
\bi{工时异质度 $V_p$}{Processing-time heterogeneity $V_p$} & 0.8 & \phpm & \ph & \ph \\
\bottomrule
\end{tabularx}

\end{table}

\begin{figure}[htbp]
\figplace{两个面板：实证 $G_Q/L^F$ 与界 $\bar B/L^F$ 随 $\bar N$（对数轴，small / main / large 档）与随 $V_p$（主实验档）的变化，标记点带运行间误差棒}{\texttt{relaxation\_diag.csv}（rounding\_gap、rounding\_bound、fluid\_lb）join \texttt{index.csv}}{opt-paper-scifig 路线 A，栏宽 88\,mm}{48mm}
\caption{实证取整差距与其理论界}\label{fig:rounding}
\end{figure}

\begin{figure}[htbp]
\figplace{左：相对最强基线增益在 $V_N$、$F$、$V_p$ 上的主效应图，带误差棒；右：双因子交互热力图}{\texttt{eval\_results.csv} 按 instance\_id join \texttt{index.csv}}{opt-paper-scifig 路线 A，双栏宽 178\,mm}{52mm}
\caption{相对最强基线增益的因子效应}\label{fig:factor}
\end{figure}

\subsection{训练与搜索行为}\label{sec:training}
本小节回答搜索与排序器训练是否稳定。图 \ref{fig:anytime} 给出固定验证集上各方法在位最优 gap 随墙钟的变化与排序器的训练曲线。

\begin{figure}[htbp]
\figplace{左：各方法验证集 gap 随墙钟的 anytime 曲线，均值 $\pm$ 标准差；右：排序器 \NumRuns{} 次训练的逐 epoch 排序损失与 top-10 recall}{\texttt{curve.csv} 与 \texttt{log.csv}}{opt-paper-scifig 路线 A，栏宽 88\,mm}{50mm}
\caption{anytime 行为与排序器训练曲线}\label{fig:anytime}
\end{figure}

\subsection{向更高多重度的泛化}\label{sec:generalization}
本小节回答方法是否只在训练多重度区间内有效。表 \ref{tab:generalization} 用冻结的排序器直接求解 large 档（$\bar N \in \{25, 50\}$）；图 \ref{fig:gen} 画出 gap 随 $\bar N$ 的变化并以阴影标出训练区间。\pending{FRES 的 gap 在 $\bar N$ 上平坦，而 job-indexed 宿主的 gap 随之增长。}

\begin{table}[htbp]
\centering\small
\caption{训练区间（$\bar N \le 12$）之外的多重度外推}\label{tab:generalization}
% generated by paper_assets/scripts/make_tables.py -- do not edit by hand

\begin{tabularx}{\linewidth}{@{}>{\hsize=1.100\hsize\raggedright\arraybackslash}X>{\hsize=1.000\hsize\raggedright\arraybackslash}X>{\hsize=0.900\hsize\raggedright\arraybackslash}X>{\hsize=1.000\hsize\raggedright\arraybackslash}X>{\hsize=0.900\hsize\raggedright\arraybackslash}X@{}}
\toprule
\textbf{\bi{方法}{Method}} & \textbf{\bi{$\bar N = 25$ gap（\%）}{$\bar N = 25$ gap (\%)}} & \textbf{\bi{$\bar N = 25$ 可行率}{$\bar N = 25$ feas.}} & \textbf{\bi{$\bar N = 50$ gap（\%）}{$\bar N = 50$ gap (\%)}} & \textbf{\bi{$\bar N = 50$ 可行率}{$\bar N = 50$ feas.}} \\
\midrule
\bi{HGA（job-indexed）}{HGA (job-indexed)} & \phpm & \ph & \phpm & \ph \\
\bi{流体进化算法（FMAE-s）}{Fluid-assisted EA (FMAE-s)} & \phpm & \ph & \phpm & \ph \\
\bi{A2 无配额压缩}{A2 no quota compression} & \phpm & \ph & \phpm & \ph \\
\bi{\textbf{FRES}}{\textbf{FRES}} & \phb & \phb & \phb & \phb \\
\bottomrule
\end{tabularx}

\end{table}

\begin{figure}[htbp]
\figplace{FRES、HGA、FMAE-s 与 A2 的 gap 随 $\bar N$ 变化，训练区间以阴影标出，$x$ 轴取对数}{\texttt{eval\_results.csv}（tier、gap）join \texttt{index.csv}（scale）}{opt-paper-scifig 路线 A，栏宽 88\,mm}{48mm}
\caption{gap 随平均多重度的变化}\label{fig:gen}
\end{figure}

\subsection{案例分析}\label{sec:case}
本小节回答调度是否合理、配额是否贴近流体分流。图 \ref{fig:case} 给出一个主实验档算例的甘特图，并列出每个类型工序的整数配额与流体比例。

\begin{figure}[htbp]
\figplace{左：FRES 在一个主实验档算例上的甘特图，同类副本同一色系，关键路径加框；右：每个类型工序的 $n_{rki}/N_r$ 与 $q^*_{rki}/N_r$ 对照柱}{单个算例的调度导出；\texttt{relaxation\_diag.csv}}{opt-paper-scifig 路线 A，双栏宽 178\,mm}{60mm}
\caption{单算例的调度与配额-流体对照}\label{fig:case}
\end{figure}

% sections/zh/6-conclusion.tex -- 总结、逐条局限、展望；与英文稿逐段对应。
\section{结论与展望}\label{sec:conclusion}
本文把多重度柔性作业车间调度问题的流体差距精确分解为取整损失与排序残差，并围绕这一恒等式构建进化搜索：维度由类型结构决定的配额染色体、输出保界且只面向残差的排序器，以及精确解码下的保底预选。\pending{在固定设计集、公开基准与外推档上，方法在同墙钟内达到不劣于各族最强基线的 gap，精确解码次数更少，取整差距的符号与理论预测一致。}

本研究的局限逐条对应设计边界。算例族固定 $R = 8$ 种类型、$m = 10$ 台机器且 $\bar N \le 50$，类型数的扩展未评测；换模、运输、故障、动态到达与批量拆分被排除（假设 A4 与 A5）；目标只有 makespan；数据为参数表分布与 Brandimarte 基准而非工厂数据，Hurink 集未纳入；未与长时限商用求解器及大语言模型自动设计的算法比较。\pending{官方代码不可得的学习型基线按复现协议重新实现，协议与结果一并说明。}

后续有三个方向。其一，同样的分解适用于任何在聚合意义上紧的松弛，可在批量与分批流变体上检验；其二，共享商图编码器可让排序器跨类型结构训练；其三，保底并集可接入其它候选源，如精确的小邻域搜索，而不改变接受准则。

\appendix
\section{术语与符号对照}
% generated by paper_assets/scripts/make_tables.py -- do not edit by hand
\begin{table*}
\centering\small
\caption{\bi{术语、译法与符号对照（全文只用一种译法）}{Terminology and notation used throughout the paper}}\label{tab:glossary}
\iflangzh

\begin{tabularx}{\linewidth}{@{}>{\hsize=1.100\hsize\raggedright\arraybackslash}X>{\hsize=1.350\hsize\raggedright\arraybackslash}Xl>{\hsize=0.550\hsize\raggedright\arraybackslash}X@{}}
\toprule
\textbf{中文术语} & \textbf{英文术语} & \textbf{缩写} & \textbf{符号} \\
\midrule
研究方案、创新点、定理、命题、猜想 & research proposal, contribution, theorem, proposition, conjecture & — & — \\
算例参数表、算例设计表、固定算例参数表 & instance parameter table, instance design table, fixed-instance parameter table & — & — \\
算法参数表、预实验、可行率、最优性间隙、独立运行 & algorithm parameter table, pilot experiment, feasibility rate, optimality gap, independent run & — & — \\
规模因子（平均多重度）与结构因子（多重度离散度、柔性度、工时异质度） & scale factor (mean multiplicity) and structural factors (multiplicity dispersion, flexibility, processing-time heterogeneity) & — & $\bar N$, $V_N$, $F$, $V_p$ \\
多重度柔性作业车间调度问题（副本数全为 1 时退化为柔性作业车间调度问题） & multiplicity flexible job-shop scheduling problem (flexible job-shop scheduling problem when every multiplicity is 1) & MFJSP, FJSP & — \\
流体-残差进化搜索（方法总名） & fluid-residual evolutionary search & FRES & — \\
配额压缩表示（个体；整数配额与优先键） & quota-compressed representation (individual; integer quota and priority key) & — & $\chi = (\mathbf{n}, \boldsymbol{\kappa})$, $n_{rki}$, $\kappa_{rk}$ \\
残差限定的商图候选排序器 & residual-scoped quotient ranker & — & $\hat G_S(a)$ \\
保底预选（三源候选数：学习器、流体启发式、随机） & safeguarded candidate preselection (per-source candidate counts) & — & $\mathcal{A}_{\mathrm{eval}}$, $K_L$, $K_H$, $K_R$ \\
流体松弛、流体变量、流体下界、机器负载对偶变量 & fluid relaxation, fluid variable, fluid lower bound, machine-load dual variable & — & $q_{rki}$, $L^F$, $\lambda_i$ \\
整数配额下界、取整界、取整损失、排序残差及其归一化值 & integer-quota lower bound, rounding bound, rounding loss, sequencing residual, normalized gap components & — & $L^Q(\mathbf{n})$, $B_i$, $\bar B$, $G_Q$, $G_S$, $g_Q$, $g_S$ \\
商图、关键块、对称感知解码器、精确解码、配额迁移 & quotient graph, critical block, symmetry-aware decoder, exact decoding, quota transfer & — & — \\
候选池、候选动作、动作变换 & candidate pool, candidate action, action transform & — & $\mathcal{A}$, $a$, $T_a$ \\
工件类型集合、机器集合与各索引 & job-type set, machine set and indices & — & $\mathcal{R}$, $M$, $r$, $c$, $i$, $o$, $o'$ \\
问题参数（副本数、工序数、候选机器集、加工时长、最短工艺路线长、big-M） & problem parameters (multiplicity, number of operations, candidate machine set, processing time, shortest route length, big-M) & — & $N_r$, $K_r$, $M_{rk}$, $p_{rki}$, $p_{oi}$, $\ell_r^{\min}$, $B$ \\
决策变量（机器分配、开工时刻、先后次序）与最大完工时间 & decision variables (machine assignment, start time, sequencing) and makespan & — & $x_{oi}$, $s_o$, $y_{oo'}$, $C_{\max}$ \\
\bottomrule
\end{tabularx}
\else

\begin{tabularx}{\linewidth}{@{}>{\hsize=1.300\hsize\raggedright\arraybackslash}Xl>{\hsize=0.700\hsize\raggedright\arraybackslash}X@{}}
\toprule
\textbf{Term} & \textbf{Abbr.} & \textbf{Symbols} \\
\midrule
research proposal, contribution, theorem, proposition, conjecture & — & — \\
instance parameter table, instance design table, fixed-instance parameter table & — & — \\
algorithm parameter table, pilot experiment, feasibility rate, optimality gap, independent run & — & — \\
scale factor (mean multiplicity) and structural factors (multiplicity dispersion, flexibility, processing-time heterogeneity) & — & $\bar N$, $V_N$, $F$, $V_p$ \\
multiplicity flexible job-shop scheduling problem (flexible job-shop scheduling problem when every multiplicity is 1) & MFJSP, FJSP & — \\
fluid-residual evolutionary search & FRES & — \\
quota-compressed representation (individual; integer quota and priority key) & — & $\chi = (\mathbf{n}, \boldsymbol{\kappa})$, $n_{rki}$, $\kappa_{rk}$ \\
residual-scoped quotient ranker & — & $\hat G_S(a)$ \\
safeguarded candidate preselection (per-source candidate counts) & — & $\mathcal{A}_{\mathrm{eval}}$, $K_L$, $K_H$, $K_R$ \\
fluid relaxation, fluid variable, fluid lower bound, machine-load dual variable & — & $q_{rki}$, $L^F$, $\lambda_i$ \\
integer-quota lower bound, rounding bound, rounding loss, sequencing residual, normalized gap components & — & $L^Q(\mathbf{n})$, $B_i$, $\bar B$, $G_Q$, $G_S$, $g_Q$, $g_S$ \\
quotient graph, critical block, symmetry-aware decoder, exact decoding, quota transfer & — & — \\
candidate pool, candidate action, action transform & — & $\mathcal{A}$, $a$, $T_a$ \\
job-type set, machine set and indices & — & $\mathcal{R}$, $M$, $r$, $c$, $i$, $o$, $o'$ \\
problem parameters (multiplicity, number of operations, candidate machine set, processing time, shortest route length, big-M) & — & $N_r$, $K_r$, $M_{rk}$, $p_{rki}$, $p_{oi}$, $\ell_r^{\min}$, $B$ \\
decision variables (machine assignment, start time, sequencing) and makespan & — & $x_{oi}$, $s_o$, $y_{oo'}$, $C_{\max}$ \\
\bottomrule
\end{tabularx}
\fi
\end{table*}

\bibliographystyle{IEEEtran}
\bibliography{refs}
\end{document}
